\centerline{\bf \Huge Числа Рамсея}

\begin{flushright}
        \scriptsize{Выживает тот, у кого есть воля к жизни.\\ Мисато Кацураги}
\end{flushright}


\problem{1} Среди любых шести человек найдутся трое попарно знакомых или трое попарно незнакомых. Докажите. Верно ли это для пяти человек?

\problem{2} Ребра клики $K_{10}$ покрашены в синий и красный цвет. Докажите, что найдется синий треугольник или красный тетраэдр.

\textit{Числом Рамсея} $R(n_1, \dots n_k)$ называется наименьшее такое $n$, что при любой раскраске ребер $n$-клики в $k$ цветов для некоторого цвета $i$ найдется $n_i-$подклика этого цвета. Заметьте, что существование чисел Рамсея мы пока не утверждаем.

\problem{3} Найдите или упростите: 
\begin{enumerate}  
    \begin{enumerate}
        \item $R(n)=$
        \item $R(1, n)=$
        \item $R(2, n)=$
    \end{enumerate} 
\end{enumerate}

\problem{4} Используя идею первых двух задач, оцените сверху:
\begin{enumerate}  
    \begin{enumerate}
        \item $R(3, 3)\leq$
        \item $R(3, 4)\leq$
        \item $R(3, 5)\leq$
        \item $R(4, 4)\leq$
    \end{enumerate} 
\end{enumerate}

\problem{5} При $a, b \geq 2$ оцените рекуррентно $R(a, b)\leq$ 

\problem{6} Оцените явно $R(a, b)\leq$

\problem{7} Представьте, что оценка в задаче $4(b)$ точна, получите противоречие. Исходя из этого пересчитайте оценки в 4 задаче. 

\problem{8} Оцените $R(3, 3, 3)$

\problem{9} При $n_1, \dots, n_k \geq 2$ оцените рекуррентно $R(n_1, \dots, n_k)$

\problem{10} Докажите, что числа Рамсея существуют!